% Part: first-order-logic
% Chapter: syntax-and-semantics
% Section: formation-sequences

\documentclass[../../../include/open-logic-section]{subfiles}

\begin{document}

\olfileid{fol}{syn}{fseq}

\olsection{دنباله‌های ساخت}

تعریف استقراییِ !!{formula}s و فن مکملِ اثبات ویژگی‌های !!{formula}s با
استقرا، رویکردی ظریف و کارآمد است. بااین‌همه، بررسی رویکردی از پایین به بالا
و گام‌به‌گام برای ساخت !!{formula}s نیز می‌تواند سودمند باشد؛ در اینجا این
کار را با استفاده از مفهوم یک \emph{دنبالهٔ ساخت} انجام می‌دهیم.
%
برای نشان‌دادن اینکه چگونه می‌توان ترم‌ها و !!{formula}s را بدین شیوه و
بی‌نیاز از ارجاع به تعریف‌های استقرایی‌شان معرفی کرد، نخست مفهوم رشته‌ای
دلخواه از نمادهای یک زبان~$\Lang L$ را معرفی می‌کنیم.

\begin{defn}[رشته‌ها]
\ollabel{defn:string}
فرض کنید $\Lang L$ زبانی مرتبهٔ اول باشد. یک \emph{رشتهٔ $\Lang L$ی}
دنباله‌ای متناهی از نمادهای~$\Lang L$ است. هرگاه زبان~$\Lang L$ به‌روشنی
از بافت معلوم باشد، اغلب رشتهٔ $\Lang L$ی را صرفاً \emph{رشته} می‌نامیم.
\end{defn}

\begin{ex}
برای هر زبان مرتبهٔ اول $\Lang L$، همهٔ !!{formula}sی $\Lang L$ی
رشته‌های $\Lang L$ی هستند، اما عکس آن برقرار نیست. برای نمونه،
\[)(\Obj v_0\lif\lexists\]
یک رشتهٔ $\Lang L$ی است، اما یک !!{formula}ی $\Lang L$ی نیست.
\end{ex}

\begin{defn}[دنباله‌های ساخت برای ترم‌ها]
\ollabel{defn:fseq-trm}
دنبالهٔ متناهیِ $\tuple{t_0,\dotsc,t_n}$ از رشته‌های $\Lang L$ی، یک
\emph{دنبالهٔ ساخت} برای ترم $t$ است هرگاه $t \ident t_n$ و برای هر
$i \leq n$، یا $t_i$ یک !!a{variable} یا !!a{constant} باشد، یا زبان
$\Lang L$ یک !!{function} $k$جایگاهی~$f$ داشته باشد و
$m_1,\dotsc,m_k < i$ چنان وجود داشته باشند که
$t_i \ident f(t_{m_1},\dotsc,t_{m_k})$. هرگاه تمایزگذاری لازم باشد،
دنباله‌های ساخت برای ترم‌ها را \emph{دنباله‌های ساخت ترم} می‌نامیم.
\end{defn}

\begin{ex}
دنبالهٔ
\[
    \tuple{\Obj c_0, \Obj v_0, \Atom{\Obj f^2_0}{\Obj c_0, \Obj v_0}, \Atom{\Obj f^1_0}{\Atom{\Obj f^2_0}{\Obj c_0, \Obj v_0}}}
\]
دنبالهٔ ساختی برای ترم $\Atom{\Obj f^1_0}{\Atom{\Obj
f^2_0}{\Obj c_0, \Obj v_0}}$ است؛ دنبالهٔ زیر نیز چنین است:
\[
    \tuple{\Obj v_0, \Obj c_0, \Atom{\Obj f^2_0}{\Obj c_0, \Obj v_0}, \Atom{\Obj f^1_0}{\Atom{\Obj f^2_0}{\Obj c_0, \Obj v_0}}}.
\]
\end{ex}

\begin{defn}[دنباله‌های ساخت برای فرمول‌ها]
\ollabel{defn:fseq-frm}
دنبالهٔ متناهیِ $\tuple{!A_0,\dotsc,!A_n}$ از رشته‌های $\Lang L$ی، یک
\emph{دنبالهٔ ساخت} برای~$!A$ است هرگاه $!A \ident !A_n$ و برای هر
$i \leq n$، یا $!A_i$ یک !!{formula} اتمی باشد یا $j,k < i$ و
!!a{variable}~$x$ چنان وجود داشته باشند که یکی از موارد زیر برقرار باشد:
\begin{enumerate}
    \tagitem{prvNot}{$!A_i \ident \lnot !A_j$.}{}%
    \tagitem{prvAnd}{$!A_i \ident (!A_j \land !A_k)$.}{}%
    \tagitem{prvOr}{$!A_i \ident (!A_j \lor !A_k)$.}{}%
    \tagitem{prvIf}{$!A_i \ident (!A_j \lif !A_k)$.}{}%
    \tagitem{prvIff}{$!A_i \ident (!A_j \liff !A_k)$.}{}%
    \tagitem{prvAll}{$!A_i \ident \lforall[x][!A_j]$.}{}%
    \tagitem{prvEx}{$!A_i \ident \lexists[x][!A_j]$.}{}%
\end{enumerate}
هرگاه تمایزگذاری لازم باشد، دنباله‌های ساخت برای فرمول‌ها را
\emph{دنباله‌های ساخت فرمول} می‌نامیم.
\end{defn}

\begin{ex}
\[
\tuple{
    \Atom{\Obj A^1_0}{\Obj v_0},
    \Atom{\Obj A^1_1}{\Obj c_1},
    (\Atom{\Obj A^1_1}{\Obj c_1} \land \Atom{\Obj A^1_0}{\Obj v_0}),
    \lexists[\Obj v_0][(\Atom{\Obj A^1_1}{\Obj c_1} \land \Atom{\Obj A^1_0}{\Obj v_0})]
}
\]
دنبالهٔ ساختی برای $\lexists[\Obj v_0][(\Atom{\Obj A^1_1}{\Obj
c_1} \land \Atom{\Obj A^1_0}{\Obj v_0})]$ است؛ دنبالهٔ زیر نیز چنین است:
\begin{multline*}
\tuple{
    \Atom{\Obj A^1_0}{\Obj v_0},
    \Atom{\Obj A^1_1}{\Obj c_1},
    (\Atom{\Obj A^1_1}{\Obj c_1} \land \Atom{\Obj A^1_0}{\Obj v_0}),
    \Atom{\Obj A^1_1}{\Obj c_1},\\
    \lforall[\Obj v_1][\Atom{\Obj A^1_0}{\Obj v_0}],
    \lexists[\Obj v_0][(\Atom{\Obj A^1_1}{\Obj c_1} \land \Atom{\Obj A^1_0}{\Obj v_0})]
}.
\end{multline*}
%
چنان‌که از مثال دوم پیداست، دنباله‌های ساخت ممکن است حاوی «اجزای زائد»
باشند: !!{formula}sی که تکراری‌اند یا در ساخت نقشی ندارند.
\end{ex}

\begin{prop}\ollabel{prop:formed}
هر !!{formula}~$!A$ در~$\Frm[L]$ دارای یک دنبالهٔ ساخت است.
\end{prop}

\begin{proof}
فرض کنید $!A$ اتمی باشد. آنگاه دنبالهٔ $\tuple{!A}$ یک دنبالهٔ ساخت
برای~$!A$ است.
%
اکنون فرض کنید $!B$ و~$!C$ به‌ترتیب دارای دنباله‌های ساخت
$\tuple{!B_0,\dotsc,!B_n}$ و $\tuple{!C_0,\dotsc,!C_m}$ باشند.
%
\begin{enumerate}
  \tagitem{prvNot}{اگر $!A \ident \lnot !B$،
      آنگاه $\tuple{!B_0,\dotsc,!B_n,\lnot !B_n}$
      یک دنبالهٔ ساخت برای~$!A$ است.}{}
  \tagitem{prvAnd}{اگر $!A \ident (!B \land !C)$،
      آنگاه $\tuple{!B_0,\dotsc,!B_n,!C_0,\dotsc,!C_m,(!B_n \land !C_m)}$
      یک دنبالهٔ ساخت برای~$!A$ است.}{}
  \tagitem{prvOr}{اگر $!A \ident (!B \lor !C)$،
      آنگاه $\tuple{!B_0,\dotsc,!B_n,!C_0,\dotsc,!C_m,(!B_n \lor !C_m)}$
      یک دنبالهٔ ساخت برای~$!A$ است.}{}
  \tagitem{prvIf}{اگر $!A \ident (!B \lif !C)$،
      آنگاه $\tuple{!B_0,\dotsc,!B_n,!C_0,\dotsc,!C_m,(!B_n \lif !C_m)}$
      یک دنبالهٔ ساخت برای~$!A$ است.}{}
  \tagitem{prvIff}{اگر $!A \ident (!B \liff !C)$،
      آنگاه $\tuple{!B_0,\dotsc,!B_n,!C_0,\dotsc,!C_m,(!B_n \liff !C_m)}$
      یک دنبالهٔ ساخت برای~$!A$ است.}{}
  \tagitem{prvAll}{اگر $!A \ident \lforall[x][!B]$،
      آنگاه $\tuple{!B_0,\dotsc,!B_n,\lforall[x][!B_n]}$
      یک دنبالهٔ ساخت برای~$!A$ است.}{}
  \tagitem{prvEx}{اگر $!A \ident \lexists[x][!B]$،
      آنگاه $\tuple{!B_0,\dotsc,!B_n,\lexists[x][!B_n]}$
      یک دنبالهٔ ساخت برای~$!A$ است.}{}
\end{enumerate}
بنا بر اصل استقرا بر !!{formula}s، هر !!{formula} دارای دنبالهٔ ساخت است.
\end{proof}

عکس این گزاره را نیز می‌توانیم اثبات کنیم. این امر مهم است، زیرا نشان می‌دهد
دو شیوهٔ ما برای تعریف فرمول‌ها هم‌ارزند: نتایج یکسانی به دست می‌دهند. همچنین
بدین معناست که می‌توانیم قضیه‌های مربوط به فرمول‌ها را با استقرای معمولی بر
طول دنباله‌های ساخت اثبات کنیم.

\begin{lem}
\ollabel{lem:fseq-init}
فرض کنید $\tuple{!A_0,\dotsc,!A_n}$ دنبالهٔ ساختی برای~$!A_n$ باشد و
$k \leq n$. آنگاه $\tuple{!A_0,\dotsc,!A_k}$ دنبالهٔ ساختی
برای~$!A_k$ است.
\end{lem}

\begin{proof}
تمرین.
\end{proof}

\begin{prob}
\olref[fol][syn][fseq]{lem:fseq-init} را اثبات کنید.
\end{prob}

\begin{thm}
\ollabel{thm:fseq-frm-equiv}
$\Frm[L]$ مجموعهٔ همهٔ رشته‌های $\Lang L$ی~$!A$ است که برای~$!A$ یک
دنبالهٔ ساخت فرمول وجود دارد.
\end{thm}

\begin{proof}
بگذارید $F$ مجموعهٔ همهٔ رشته‌های نمادهای زبان~$\Lang L$ باشد که دنبالهٔ
ساخت دارند. در \olref[fol][syn][fseq]{prop:formed} دیدیم که
$\Frm[L] \subseteq F$؛ پس اکنون شمول عکس را اثبات می‌کنیم.

فرض کنید $!A$ دارای دنبالهٔ ساخت $\tuple{!A_0,\dotsc,!A_n}$ باشد.
با استقرای قوی بر~$n$ اثبات می‌کنیم که $!A \in \Frm[L]$.
فرض استقرای ما این است که هر رشته از نمادها که دنبالهٔ ساختی با شاخص پایانی
$m < n$ دارد، در $\Frm[L]$ است.
بنا بر تعریف دنبالهٔ ساخت، یا $!A_n$ اتمی است یا باید $j,k < n$ و
!!a{variable}~$x$ وجود داشته باشند، به‌گونه‌ای که یکی از موارد زیر برقرار
باشد:
\begin{enumerate}
    \tagitem{prvNot}{$!A \ident \lnot !A_j$.}{}%
    \tagitem{prvAnd}{$!A \ident (!A_j \land !A_k)$.}{}%
    \tagitem{prvOr}{$!A \ident (!A_j \lor !A_k)$.}{}%
    \tagitem{prvIf}{$!A \ident (!A_j \lif !A_k)$.}{}%
    \tagitem{prvIff}{$!A \ident (!A_j \liff !A_k)$.}{}%
    \tagitem{prvAll}{$!A \ident \lforall[x][!A_j]$.}{}%
    \tagitem{prvEx}{$!A \ident \lexists[x][!A_j]$.}{}%
\end{enumerate}
اکنون به تفکیک حالت‌ها استدلال می‌کنیم. اگر $!A_n$ اتمی باشد، آنگاه
$!A_n \in \Frm[L]$. حال فرض کنید
$!A \ident (!A_j \land !A_k)$. بنا بر
\olref[fol][syn][fseq]{lem:fseq-init}،
$\tuple{!A_0,\dotsc,!A_j}$ و $\tuple{!A_0,\dotsc,!A_k}$ به‌ترتیب
دنباله‌های ساختی برای $!A_j$ و~$!A_k$ هستند. چون هر یک قطعهٔ آغازینِ سره‌ای
از دنبالهٔ ساخت~$!A$ است، شاخص پایانیِ هر دو کوچک‌تر از~$n$ است. بنابراین،
بنا بر فرض استقرا، $!A_j$ و~$!A_k$ در~$\Frm[L]$ هستند و بنا بر تعریف
!!a{formula}، $(!A_j \land !A_k)$ نیز چنین است. حالت‌های دیگر با
استدلالی موازی به دست می‌آیند.
\end{proof}

دنباله‌های ساخت ترم ویژگی‌هایی مشابه دنباله‌های ساخت !!{formula}s دارند.

\begin{prop}
\ollabel{prop:fseq-trm-equiv}
$\Trm[L]$ مجموعهٔ همهٔ رشته‌های $\Lang L$یِ $t$ است که برای~$t$ یک
دنبالهٔ ساخت ترم وجود دارد.
\end{prop}

\begin{proof}
تمرین.
\end{proof}

\begin{prob}
\olref[fol][syn][fseq]{prop:fseq-trm-equiv} را اثبات کنید.
راهنمایی: راهبردی مشابه راهبرد به‌کاررفته در اثبات
\olref[fol][syn][fseq]{thm:fseq-frm-equiv} به کار ببرید.
\end{prob}

دو نوع «جزء زائد» می‌تواند در دنباله‌های ساخت ظاهر شود: عناصر تکراری و
عناصری که به ساخت فرمول یا ترم مربوط نیستند. با بررسی دنباله‌های ساخت کمینه
می‌توانیم هر دو نوع را حذف کنیم.

\begin{defn}[دنباله‌های ساخت کمینه]
\ollabel{defn:minimal-fseq}
دنبالهٔ ساخت $\tuple{!A_0, \ldots, !A_n}$ برای فرمول~$!A$ یک
\emph{دنبالهٔ ساخت کمینه} برای~$!A$ است هرگاه برای هر دنبالهٔ ساخت
دیگر~$s$ برای~$!A$، طول~$s$ بزرگ‌تر یا مساوی~$n+1$ باشد.

به همین ترتیب، دنبالهٔ ساخت $\tuple{t_0, \ldots, t_n}$ برای ترم~$t$ یک
\emph{دنبالهٔ ساخت کمینه} برای~$t$ است هرگاه برای هر دنبالهٔ ساخت
دیگر~$s$ برای~$t$، طول~$s$ بزرگ‌تر یا مساوی~$n+1$ باشد.
\end{defn}

توجه کنید که یک فرمول یا ترم می‌تواند بیش از یک دنبالهٔ ساخت کمینه داشته
باشد، اما این دنباله‌ها دقیقاً رشته‌های یکسانی را در بر خواهند داشت.

\begin{prop}
\ollabel{prop:subformula-equivs}
گزاره‌های زیر هم‌ارزند:
\begin{enumerate}
    \item $!B$ یک !!{subformula} از~$!A$ است.
    \item $!B$ در هر دنبالهٔ ساخت~$!A$ رخ می‌دهد.
    \item $!B$ در یک دنبالهٔ ساخت کمینهٔ~$!A$ رخ می‌دهد.
\end{enumerate}
\end{prop}

\begin{proof}
تمرین.
\end{proof}

\begin{prob}
\olref[fol][syn][fseq]{prop:subformula-equivs} را اثبات کنید.
\end{prob}

\begin{history}
دنباله‌های ساخت را ریموند اسمولیان در کتاب درسی \emph{منطق مرتبهٔ اول}
\citep{Smullyan1968} معرفی کرد. ویژگی‌های دیگری از دنباله‌های ساخت را
\citet{Zuckerman1973} اثبات کرد.
\end{history}

\end{document}
